\section{Conclusion}
\label{sec:conclusion}

This program set out to ask whether the harmonized baseline of three diagnostic cohorts --- bipolar
disorder, schizophrenia, and depression --- can be reorganized into a transdiagnostic, biology-aware
coordinate system that \emph{describes}, \emph{tracks}, \emph{forecasts}, and ultimately \emph{guides}
care better than the diagnostic labels themselves. It pursued that question under a discipline that
refuses to collapse its four layers (cohorts $\to$ dimensions $\to$ strata $\to$ prognosis/treatment) or
to manufacture certainty it has not earned: an observed-data likelihood with \emph{no imputation},
diagnosis held as metadata rather than a feature, dimensions fixed on the baseline visit and only ever
\emph{validated} forward, and every estimate carrying its own uncertainty into the next.

The report answers the question in escalating senses. The map \textbf{exists}: one global
Bayesian bifactor/ESEM on observed cells yields a nine-dimension structure --- a general
functional-burden factor $G$ and eight weakly-correlated specific axes --- in which biology (metabolic,
inflammatory load) is the \emph{least} severity-entangled domain (metabolic$\sim$$G$ $+0.12$,
inflammatory$\sim$$G$ $+0.07$, against $0.39$ for cognition and $0.42$ for sleep), riding on axes that
clinical severity does not see. It \textbf{organizes}: on those coordinates patients form not discrete
biotypes but a graded \emph{continuum} --- a single-Gaussian falsification null cannot be separated from
the real data (best silhouette $0.140$ vs.\ $0.140\pm0.003$) --- across which four stable soft archetypes
mark quasi-independent corners of biology, symptoms, and severity, cutting across the DSM-5 subtypes
(adjusted Rand $\approx 0$) while describing the data more tightly than diagnosis does. It
\textbf{persists}: scored forward onto follow-up without re-estimation, the measurement is invariant (all
five backbone axes $\varphi\geq0.95$) and the geometry replays --- durable biology (metabolic test--retest
$\text{ICC}=0.91$), moving symptoms. It \textbf{predicts}: a baseline coordinate forecasts two-year
functioning incrementally beyond diagnosis, severity, and the present outcome
($\Delta\text{ELPD}=+59$ for the archetypes), with the biological corner carrying the worst functional
prognosis ($27\%$ remission, against $60\%$ at the low-burden pole). And it does \textbf{not}, on the
evidence available, \textbf{prescribe}: with treatment exposures recovered from the per-cohort source
dictionaries and run through a proper overlap-gated causal pipeline, the map does not reliably
\emph{moderate} treatment-as-usual response.

The honest reading separates two bars the field too often conflates. On \textbf{scientific validity} the
program succeeds: the map is real, identifiable, stable, transdiagnostic, a continuum rather than a set of
biotypes, with biology the least severity-entangled domain and a genuine, robustness-surviving prognostic
signal for functioning that is co-informative with DSM-5 and even \emph{survives} adjustment for the
treatments patients received. On \textbf{clinical utility, in the strong sense}, it does not: the
individual-level prognostic increment is small (the $+59$-ELPD group-level gain collapses to a
functional-remission $\Delta\text{AUC}$ of only $+0.011$, $0.745\to0.756$), and the map does not moderate
treatment in observational data. Several of the program's most secure findings are therefore
\emph{negative or modest} --- no biotypes, no reliable treatment moderation, a group-level rather than
individual prognostic gain --- and they are presented as the contribution, not buried beneath the
seductive numbers.

The result must also be carried with its tensions stated, not smoothed over, because each marks the exact
edge of what the data support. The predictive object is a \emph{phenotype}, not an isolated axis: the
forecast lives in the fuller four-archetype representation --- biology read as a corner of the whole space
--- and not in a lone three-axis ``durable block'', which on its own is only ambiguous. The prognostic
gain is real but \emph{group-level}, sharpening who on average is heading where rather than calling an
individual remission. It is \emph{course-dependent}, carried by the episodic bipolar and depression
patients and near-null in the baseline-saturated schizophrenia cohort. The tessellation has \emph{no
privileged number of regions} --- an operationally awkward but honest consequence of the continuum, with
the ``operative'' resolution deferred to, and then dissolved by, the prognosis question (the
continuous/archetype encoding dominates any hard partition). Durability is cleanest in the
\emph{variance} decomposition (the error-corrected trait/state ICCs); the parallel geometry route is
precision-confounded and is treated as corroboration, not as the primary signal. And the map's thinnest
axis, \emph{mania} (two indicators), is reliable-looking but prognostically and temporally uninformative,
and is flagged as such rather than over-read.

What the work establishes, then, is a transdiagnostic biological map that \textbf{complements} the
diagnostic system rather than replacing it: co-informative with DSM-5 prognostically, a tighter
description of the coordinates than diagnosis, and a continuous, uncertainty-aware representation of where
a patient sits in biology-and-severity space. What it does not establish is an individual risk calculator
or a prescribing rule. In a literature replete with biological-validity evidence advanced as though it
were clinical utility, a study that measures \emph{both} and reports the modest and null results as such
is itself a contribution --- a calibrated counterweight, and a reusable, principled pipeline (no
imputation, uncertainty propagated end-to-end, fixed-then-validated objects, overlap-gated causal
inference) for asking these questions without over-claiming.

The boundary is drawn by the data, not by the method, and the work names its own exits
(Chapter~\ref{sec:limitations}). Individual-level clinical utility would require incident-event outcomes
rather than re-administered scales; treatment \emph{selection} would require randomized or trial-arm data
rather than treatment-as-usual --- a defined future study (M5b); and any clinical claim would require
external validation beyond this internally-validated baseline cohort. Each is a defined next study, and the
present program is the methods-and-feasibility foundation on which they can be built.

\begin{keyresult}[The deliverable, stated plainly]
A stable, transdiagnostic biological map that is a \emph{continuum} rather than a set of
biotypes, whose durable biological corner forecasts two-year functioning beyond diagnosis and severity at
the group level, and that --- on observational data --- does not yet guide treatment. A trustworthy map,
and an honest account of how far it reaches.
\end{keyresult}
